% Part: history
% Chapter: biographies 
% Section: emmy-noether

\documentclass[../../../include/open-logic-section]{subfiles}

\begin{document}

\olfileid{his}{bio}{noe}

\olsection{ایمی نوئتر}

\olphoto{noether-emmy}{ایمی نوئتر}

ایمی نوئتر (\textsc{نَر} ٹر) 23 مارچ 1882 نوں جرمنی دے ایرلانگن شہر وچ پڑھے لکھے اُچلے درمیانے طبقے دے ٹبّر وچ جمی۔ اوہنوں «اج دی الجبرا دی ماں» آکھیا جاندا اے۔ عورتاں دی تعلیم اگے وڈیاں رکاوٹاں دے باوجود نوئتر نے ریاضی تے طبیعیات دوہاں وچ نئی راہ کھولھن والا کم کیتا۔ اوس ویلے جرمنی وچ کُڑیاں نوں فن پڑھائے جاندے سن تے اوہناں نوں کالج دی تیاری والے سکولاں وچ جان دی کھُل نہیں سی۔ پر گوٹنگن تے ایرلانگن دیاں یونیورسٹیاں وچ داخلے توں بناں جماعتاں سُنن پچھوں (ایرلانگن وچ اوہدا پیو ریاضی دا پروفیسر سی)، نوئتر آخر 1904 وچ ایرلانگن وچ طالب علم وجوں داخل ہو سکی، جدوں اوتھوں دے قاعدے بدل کے عورتاں نوں داخلے دی کھُل دتی گئی۔ اوہنے 1907 وچ ریاضی وچ ڈاکٹریٹ لئی۔

اپنی قابلیت دے باوجود نوئتر نوں اپنے پیشے وچ بہت مخالفت دا سامنا کرنا پیا۔ 1908 توں 1915 تک اوہنے ایرلانگن وچ تنخواہ توں بناں پڑھایا۔ ایس سمیں اوہ ڈیوڈ ہلبرٹ دی نظر چڑھی، جیہڑا اوس ویلے دنیا دے وڈے ریاضی داناں وچوں اک سی۔ ہلبرٹ نے اوہنوں گوٹنگن سَدیا۔ پر عورتاں نوں پروفیسری لین دی کھُل نہیں سی۔ نوئتر بس ہلبرٹ دے ناں تھلّے، فیر تنخواہ توں بناں، لیکچر دے سکدی سی۔ ایس سمیں اوہنے اوہ گل ثابت کیتی جیہنوں ہُن نوئتر دا قضیہ آکھدے نیں، تے جیہڑا اج وی نظری طبیعیات وچ ورتیا جاندا اے۔ آخر 1919 وچ نوئتر نوں پڑھاؤن دا حق مل گیا۔ دسیا جاندا اے کہ یونیورسٹی دے ساتھی جدوں ہن وی مخالفت کردے رہے تاں ہلبرٹ نے جواب دِتا، «حضرات، اساتذہ دی سینیٹ کوئی حمام نہیں۔»

1920 دی دہائی دے پچھلے سالاں وچ اوہنے تجریدی الجبرا اُتے دھیان دِتا، تے اوہدے کم نے ایس شعبے نوں بدل کے رکھ دِتا۔ اپنے ثبوتاں وچ اوہ اکثر چڑھدی زنجیر دی شرط ورتدی سی، جیہڑی دسدی اے کہ کجھ خاص سیٹاں دی سختی نال ودھدی کوئی لامتناہی زنجیر نہیں۔ مثلاً ہُن نوئتری حلقے آکھیاں جاندیاں کجھ الجبرائی ساختاں دی خاصیت اے کہ اوہناں وچ آئیڈیلاں دی کوئی لامتناہی لڑی $I_1
\subsetneq I_2 \subsetneq \dots$ نہیں۔ ایہ شرط کسے وی جزوی ترتیب تک ودھائی جا سکدی اے (الجبرا وچ ایہ ذیلی سیٹ دے تعلق نال ترتیب دتے آئیڈیلاں دی خاص صورت بارے اے)۔ اسیں ایس دی دوہری، لہندی زنجیر دی شرط وی ویکھ سکدے آں، جتھے جزوی ترتیب وچ سختی نال \emph{گھٹدی} ہر لڑی آخر مُک جاندی اے۔ جے کوئی جزوی ترتیب لہندی زنجیر دی شرط اُتے پوری اُترے، تاں اوس ترتیب اُتے استقرا اویویں ورت سکدے آں جیویں $<$ ترتیب اُتے~$\Nat$ وچ ورتدے آں۔ ایہو جیہیاں ترتیباں نوں \emph{بنیاد والیاں} یا \emph{نوئتری} آکھدے نیں، تے اوہناں دے ثبوت دے اصول نوں \emph{نوئتری استقرا} آکھدے نیں۔

نوئتر یہودی سی۔ جدوں 1933 وچ نازی راج وچ آئے تاں اوہنوں عہدے توں کڈھ دِتا گیا۔ چنگے بھاگیں نوئتر امریکہ ہجرت کر سکی، جتھے برِن مار، پنسلوانیا وچ اوہنوں عارضی عہدہ مل گیا۔ اوتھے رہندیاں اوہنے پرنسٹن وچ وی لیکچر دتے، بھاویں اوہنوں لگدا سی کہ یونیورسٹی عورتاں نوں جی آیاں نہیں آکھدی \citep[81]{Dick1981}۔ 1935 وچ بچہ دانی دی رسولی کڈھن لئی نوئتر دا آپریشن ہویا۔ اوہ آپریشن توں ہوئی انفیکشن نال مر گئی تے برِن مار وچ دفنائی گئی۔

\begin{reading}
نوئتر دی سوانح عمری لئی \citet{Dick1981} ویکھو۔ پیریمیٹر انسٹیٹیوٹ فار تھیوریٹیکل فزکس دے نوئتر دی زندگی تے اثر بارے لیکچر انٹرنیٹ اُتے ملدے نیں \citep{Perimeter2015}۔ جے تسی پڑھ پڑھ کے تھک گئے او، تاں \emph{Stuff You Missed in History Class} دا نوئتر دی زندگی تے اثر بارے اک پوڈکاسٹ اے \citep{Frey2015}۔ نوئتر دیاں اکٹھیاں لکھتاں اصل جرمن وچ ملدیاں نیں \citep{Noether1983}۔
\end{reading}

\end{document}
